\documentclass[11pt, oneside]{article}   	% use "amsart" instead of "article" for AMSLaTeX format
\usepackage{geometry}                		% See geometry.pdf to learn the layout options. There are lots.
\geometry{letterpaper}                   		% ... or a4paper or a5paper or ... 

\usepackage[parfill]{parskip}    		% Activate to begin paragraphs with an empty line rather than an indent
\usepackage{graphicx}				% Use pdf, png, jpg, or eps§ with pdflatex; use eps in DVI mode
								% TeX will automatically convert eps --> pdf in pdflatex		
\usepackage{amssymb}
%\usepackage{cite}
\usepackage[
backend=biber,
style=verbose,
%style=nature,
sorting=nyt
]{biblatex}
%\addbibresource{../Chemistry.bib}
\bibliography{../Chemistry}

% numbered examples
\usepackage{gb4e}
\usepackage{enumitem}
\usepackage{cancel}
\usepackage{amsmath}

\usepackage[version=4]{mhchem}
\usepackage{lewis}

\usepackage{url}
\usepackage{hyperref}


% Scientific Laws
\newtheorem{definition}{Definition}
\newtheorem{law}{Law}
\newtheorem{theory}{Theory}
\newtheorem{hint}{Hint}

\title{Module 15: Chemical Equilibrium}
\author{Mark Peever\\ \texttt{mpeever@gmail.com}}
\date{April 17 -- 24, 2026}

\begin{document}
\maketitle

\section{Overview}
\begin{enumerate}
\item chemical reactions can occur both ``forwards" and ``backwards"
\item eventually reactions settle into an \textbf{Equilibrium} where the ``forward" and ``backward" reactions are in a balance
\item we can describe the balance mathematically using the concept of \textbf{Reaction Rate} from Module 14
\item \textbf{Le Chatelier's Principle} tells us how \textbf{stresses} will affect Chemical Equilibrium
\item \textbf{pH} is a way to quantify Chemical Equilibrium in Acid/Base reactions
\end{enumerate}

\section{Chemical Equilibrium}
\begin{definition}[Chemical Equilibrium]\label{defn:chemical-equilibrium}
The point at which both the forward and reverse reactions in a chemical equation have equal reaction rates:
When this occurs the amount of each substance in the chemical reaction will not change, despite the fact that both reactions still proceed.
\end{definition}

\begin{itemize}
\item it turns out chemical reactions can happen both ``forward" and ``backward":
\begin{itemize}
\item \ce{NH3 (aq) + H2O (l) -> NH4+ (aq) + OH- (aq)}
\item \ce{NH3 (aq) + H2O (l) <- NH4+ (aq) + OH- (aq)}
\end{itemize}
\item we combine both reactions into a single ``bi-directional" reaction:
\begin{itemize}
\item  \ce{NH3 (aq) + H2O (l) <=> NH4+ (aq) + OH- (aq)}
\end{itemize}
\item when we first combine the reactants (in this case \ce{NH3} and  \ce{H2O}), the reaction will run ``forwards"
\item but it will eventually\footnote{By ``eventually" I mean ``after some non-zero period of time."}
begin to run ``backwards" too
\item at some point the left-to-right and right-to-left reactions will be occurring at the same rate, so that we'll have a constant amount of each reactant and/or product
\item when the reaction rates stabilize, we say our reaction as reached \textbf{Chemical Equilibrium} (see Definition \ref{defn:chemical-equilibrium}, page \pageref{defn:chemical-equilibrium}) 
\item reactions rates \emph{in both directions} are driven by \emph{concentrations} of the reactants, just like in Module 14
\item it takes time for a reaction\footnote{Really, a \emph{pair} of reactions.} to achieve equilibrium
\item once a reaction reaches equilibrium, it will tend to stay there until some external ``stress" acts on it
\footnote{Whoa! This sounds almost exactly like Newton's Second, doesn't it?}
\end{itemize}

\section{The Equilibrium Constant}

\begin{definition}[Equilibrium Constant]\label{defn:equilibrium-constant}
For some Chemical Reaction:
 
 \begin{center}
 \ce{aA + bB <=> cC + dD}
 \end{center}
 
the Equilibrium Constant is calculated as:\\
$$K = \frac{[C]^{c}_{eq} \times [D]^{d}_{eq}} {[A]^{a}_{eq} \times [B]^{b}_{eq}}$$
\end{definition}



\begin{itemize}
\item we quantify Chemical Equilibrium with an \emph{Equilibrium Constant} (see Defintion \ref{defn:equilibrium-constant}, page \pageref{defn:equilibrium-constant}))
\item notice this constant will be inverted if we write the equations right-to-left instead of left-to-right
\item $K$ here will depend on the specific reaction \emph{and} the environment
\footnote{For example, altering temperature will alter $K$, but don't let's get into that quite yet\ldots}
\item notice that the exponents on the concentrations are the \emph{stoichiometric coefficients}, so this isn't quite the same thing we saw in Module 14
\item we can make some helpful generalizations about a chemical reaction just looking at the equilibrium constant (see Hint \ref{hint:equilibrium-constant-generaiizations}, page \pageref{hint:equilibrium-constant-generaiizations}):
\begin{itemize}
\item when $K$ is large, the equilibrium is weighted to the products (the right-hand-side)
\item when $K$ is small, the equilibrium is weighted to the reactants (the left-hand-side)
\item when $K$ is near $1$, the equilibrium is balanced between the two sides
\end{itemize}
\item again, notice if we write the equation ``backwards" we'll \emph{flip} the Equilibrium Constant
\item our Chemical Equilibrium Constant is only for concentrations, \emph{i.e.} it's only for chemicals in solution:
\begin{itemize}
\item it doesn't apply to solids
\item it doesn't apply to liquids\footnote{But \emph{do} include \emph{aqueous} substances!}
\end{itemize}
\end{itemize}

\begin{hint}[Equilibrium Constant Generalizations]\label{hint:equilibrium-constant-generaiizations}
We can generalize that:\\
\begin{itemize}
\item when $K$ is large, the equilibrium is weighted to the products (the right-hand-side)
\item when $K$ is small, the equilibrium is weighted to the reactants (the left-hand-side)
\item when $K$ is near $1$, the equilibrium is balanced between the two sides
\end{itemize}
\end{hint}

\section{What the Chemical Equilibrium Constant Tells Us}
\begin{itemize}
\item we can use an \emph{Equilibrium Constant} (see Defintion \ref{defn:equilibrium-constant}, page \pageref{defn:equilibrium-constant}) to determine whether a reaction is at equilibrium
\item we can also predict where the equilibrium will be
\end{itemize}

\section{Le Chatelier's Principle}

\begin{definition}[Le Chatelier's Principle]\label{defn:le-chatelier}
When a \emph{stress} (such as a change in concentration, pressure, or temperature) is applied to an equilibrium,
the reaction will shift in a way that relieves the stress and restores equilibrium.
\end{definition}

\begin{itemize}
\item once chemical equilibrium is reached, a reaction will tend to stay at equilibrium indefinitely
\item if we alter the reaction by changing something (\emph{e.g.} we add more of a given substance, or maybe take some away), then it will adjust to a new equilibrium
\item for example:
\begin{enumerate}
\item if we start with our favorite reaction:\\ \begin{center}\ce{NH3 (aq) + H2O (l) <=> NH4+ (aq) + OH- (aq)}\end{center}
\item and if we pour in more \ce{NH3}
\item then our reaction will shift \emph{right} to use up the additional \ce{NH3}
\item it will eventually settle into a new equilibrium
\end{enumerate}
\item we can use Le Chatelier to force reactions to run ``to completion" by adding stresses to reactions at equilibrium
\item \emph{e.g.} we could remove \ce{OH-} from our equilibrium so that more \ce{NH3} is used by ``trying" to replace it
\item Le Chatelier treats enthalpy and entropy as stresses, see Hint \ref{hint:le-chatelier-pressure} and Hint \ref{hint:le-chatelier-temperature}
\end{itemize}

\begin{hint}\label{hint:le-chatelier}
Le Chatelier's Principle ignores solids and liquids as a source of stress to the equilibrium.
\end{hint}

\begin{hint}\label{hint:le-chatelier-pressure}
When an equilibrium is subjected to an increase in pressure, it will shift away from the side with the largest number of gas molecules.
If pressure decreases, the equilibrium will shift toward the side with the largest number of gas molecules.
If there are no gases in the equation or if the number of gas molecules is equal on both sides, nothing will happen.
\end{hint}

\begin{hint}\label{hint:le-chatelier-temperature}
When temperature is raised, an equilibrium will shift away from the side of the equation that contains energy.
When temperature is lowered, an equilibrium will shift toward the side of the equation that contains energy.
\end{hint}

\section{Acid/Base Equilibria}

\begin{definition}[Acid Ionization Reaction]\label{def:acid-ionization}
The reaction in which an \ce{H+} separates from an acid molecule so that it can be donated in another reaction.
\end{definition}

\begin{definition}[Acid Ionization Constant]\label{def:acid-ionization-constant}
The equilibrium constant for an acid's ionization reaction.
\end{definition}

\begin{definition}[Base Ionization Constant]\label{def:base-ionization-constant}
The equilibrium constant for an base's ionization reaction.
\end{definition}

\begin{itemize}
\item remember that an acids and bases \emph{ionize} so in aqueous solution:
\begin{itemize}
\item \ce{HCl (aq) + H2O (l) <=> H+ (aq) + Cl-}(see Definition \ref{def:acid-ionization}, p. \pageref{def:acid-ionization})
\footnote{We can use \ce{H3O+} and \ce{H+ (aq)} interchangeably (\emph{why?}), so we'll follow the text here.}
\item \ce{NH3 (aq) + H2O (l) <=> NH4+ (aq) + OH- (aq)}
\end{itemize}
\item for acids, we refer to $K_{a}$, which is a special case of $K$
\item $K_{a}$ is the \emph{Acid Ionization Constant}, see Definition \ref{def:acid-ionization-constant}, page \pageref{def:acid-ionization-constant}
\item well, it turns out that ionization reactions can be classified as \emph{strong} or \emph{weak} based on their $K_{a}$ values
\item strong acids include:
\begin{itemize}
\item \ce{HCl}
\item \ce{H2SO4}
\item \ce{HNO3}
\end{itemize}
\item weak acids include:
\begin{itemize}
\item \ce{C2H4O2} (acetic acid)
\item \ce{C6H8O6} (ascorbic acid, Vitamin C)
\item \ce{H3PO4} (phosphoric acid)
\end{itemize} 
\item we can bases the same way, with $K_{b}$ (see Definition \ref{def:base-ionization-constant}, page \pageref{def:base-ionization-constant})
\end{itemize}

\section{The pH Scale}

\begin{itemize}
\item we can measure the acidity or alkalinity\footnote{``Alkaline" is the older term for ``basic."} of a solution with the \emph{pH} scale
\footnote{\emph{pH} means ``potential hydrogen"}
\item we calculate pH as: $pH = - \log_{10} [H^{+}]$
\item pH is measured on a scale from $0$ --- $14$
\item higher numbers are more basic
\item lower numbers are more acidic
\item pure water has a pH of $7$
\end{itemize}





\section{Homework}
Review Problems: p. 522 \# 1--10 (not to be turned in)\\
Practice Problems: pp. 523--524 \# 1--10 (due 2026-05-01)\\
 
\nocite{feynman-entropy-lecture-video} 
\nocite{reading-feynman-entropy}
\nocite{chem:libretexts:energy-diagrams}
\nocite{wile-chem-2}
%\bibliography{../Chemistry}
%\bibliographystyle{apalike}

\printbibliography

\clearpage
\section*{Answers to Practice Problems}
Answers to Practice Problems pp. 523--524 \# 1--10:
\begin{enumerate}
\item $2.0 M^{2}$ 
\item $4.7 \times 10^{-5} M^{5}$\\ (Hint: how do solids affect $K$?)

\item 
\begin{enumerate}
\item \ce{2NO2 (g) + Cl2 (g) -> 2NOCl (g) + O2 (g)}
\item \ce{2SO3 (g) <=> 2SO2 (g) + O2 (g)}
\item \ce{2NO2 -> N2O4}
\end{enumerate}

\item No, the reaction is not at equilibrium and will shift \emph{right}.
\item Yes, the reaction is at equilibrium

\item 
\begin{enumerate}
\item There is no effect on the equilibrium\\ (Hint: what does Le Chatelier's Principle say about solids?)
\item There is no effect on the equilibrium\\ (Hint: what does Le Chatelier's Principle say about liquids?)
\item The equilibrium will shift \emph{right} to produce more \ce{CaO}.
\end{enumerate}

\item 
\begin{enumerate}
\item The equilibrium will shift \emph{right} and will lower \ce{[H2]}.
\item The equilibrium will shift \emph{left} and will lower \ce{[HF]}.
\end{enumerate}

\item 
\begin{enumerate}
\item The equilibrium will shift \emph{left} and will raise \ce{[H2]}.
\item The equilibrium will shift \emph{right} and will raise \ce{[NH3]}.
\item The equilibrium will shift \emph{left} and will raise \ce{[N2]}.
\end{enumerate}

\item $K_{a} = \frac{ [ H^{+} ] [CN^{-}] }{ [HCN]  }$
\item $K_{b} =  \frac{ [ HSO_{4}^{-} ]  [OH^{-}] } { [SO_{4}^{2-}] }$
\end{enumerate}
\end{document}  